\documentclass[12pt]{article}
\usepackage[margin=1in]{geometry}
\usepackage{enumitem}
\usepackage{titlesec}
\usepackage{parskip}
\usepackage{graphicx}
\usepackage{hyperref}
\usepackage{fontenc}

\title{\textbf{Cook and Levi 1990 RG Notes}\\
\large Cook, Karen Schweers, and Margaret Levi, eds. \textit{The Limits of Rationality}.\\
Chicago: University of Chicago Press, 1990.}
\author{}
\date{}

\begin{document}

\maketitle

\section{Main Idea}



%--------------------------------------------------
\section*{Overarching Notes}
%--------------------------------------------------

\begin{itemize}[leftmargin=*, itemsep=4pt]
    \item Overarching: what are the limits of rational choice theory, and what is its future? Claim is that the future of rational choice analysis is based in studies of institutions and norms
    \item Shares the critique/belief that rational choice models are often oversimplified in social science environments where economic principles (price theory, competition, etc.) are not relevant. Extension of game theory to focus on strategic interactions has done some to alleviate this by emphasizing the importance of group decisions, but has not been sufficiently defeated
\end{itemize}


\section{Elster: ``When Rationality Fails''}

\begin{itemize}
    \item ``Failure to recognize the indeterminacy of rational choice theory can lead to irrational behavior.''
    \item RCT does not identify what our goals should be, but does identify what should be done in order to achieve some specified goal(s)
    \item Under rationality, beliefs must be optimally related to the evidence available to the agent (kind of quote)
    \item Veyne 1976: ``To serve his master well, a slave must have some independence of execution: beliefs born of passion serve passion badly''
    \item Claim: some goals and utilities are human and cannot be negated even in the search for greater utility (for example, the desire for freedom cannot be broadly eliminated)
    \item Hyperrationality: ``the failure to recognize the failure of rational choice theory to yield unique prescriptions or predictions.''
    \item ``Rationality is presupposed by any competing theory of motivation, whereas rationality itself does not presuppose anything else. On grounds of parsimony, therefore, we should begin by assuming nothing but rationality.''
    \item Social norms can be seemingly irrational, but it is perhaps following them which produces rationality -- norms may be the strongest alternative to RCT if you assume that it does not itself result in rationality
\end{itemize}


\section{Tversky \& Kahneman: ``Rational Choice and the Framing of Decisions"}

\begin{itemize}
    \item Deviations from predicted rational outcomes are significant, potentially disqualifyingly so; concludes that the normative and descriptive analyses ``cannot be reconciled''
    \item Von Neumann Morgenstern utility: cancellation, transitivity, dominance, and invariance
\end{itemize}


\section{Machina: ``Choice Under Uncertainty: Problems Solved \& Unsolved''}

\begin{itemize}
    \item Questions arising about theories of choice under uncertainty, which has come to face a handful of challenges
\end{itemize}


\section{Stigler \& Becker: ``De Gustibus Non Est Disputandum''}

\begin{itemize}
    \item Foundational rooting in tastes as an exogenous and generally-immovable characteristic
    \item ``We are proposing the hypothesis that widespread and/or persistent human behavior can be explained by a generalized calculus of utility-maximizing behavior, without introducing the qualification `tastes remaining the same.'''
    \item Rational choice as the best available tool to explain behavior without making assumptions that tastes have changed -- that is, that allows for (but does not require) consistent tastes
    \item Fashion as a counterexample; ``The obvious method of reconciling fashion with our thesis is to resort again to the now familiar argument that people consume commodities and only indirectly do they consume market goods, so fashions in market goods are compatible with stability in the utility function of commodities.''
    \item General argument is that you can reconcile changing behavior with static tastes
\end{itemize}


\section{Taylor: ``Cooperation and Rationality: Notes on the Collective Action Problem and Its Solutions''}


\begin{itemize}
    \item Elster's definition of collective action problem:
    \begin{itemize}
        \item Strong: Prisoners' Dilemma
        \item Weak: (1) everyone prefers universal cooperation to universal noncooperation, and (2) cooperation is ``individually unstable and individually inaccessible.''
    \end{itemize}
    \item Two types of collective action solutions:
    \begin{enumerate}
        \item Spontaneous / Internal: no changes to the game, preferences, or beliefs
        \item External: changing the game, preferences, or beliefs
    \end{enumerate}
    \item External solutions can be either decentralized, such that the initiative or responsibility for enacting the solution is broadly shared, or centralized, such that the solution authority is concentrated in the hands of some subset of the members (assume some sort of domination outcome here?). Centralized is taken to be consubstantial with the state
    \item Highlights the importance of repetition, group norms, and an ability to enforce said norms as a core component of solving collective action problems
    \item One external solution is the political entrepreneur: coordination mechanism. Must change beliefs or preferences (usually occurs by adding knowledge)
    \item Property rights as a possible solution to collective action problems: if the impacts of consumption can be totally internalized within the individual, then such problems may be avoided entirely
    \item All seems very similar to Ostrom's operationalization of the Folk theorem
\end{itemize}


\section{Coleman: ``Norm-Generating Structures''}


\begin{itemize}
    \item Provides a comprehensive (but rather esoteric) definition of norms
    \item ``The demand for a norm will arise when an action by one actor imposes externalities on other actors. When the externalities are positive, the demand is for a prescriptive norm, to encourage or induce the focal action. When the externalities are negative, the demand is for a proscriptive norm, to discourage the focal action.''
    \item Norms emerge as behaviors within (and surrounding) repeated PD interactions. Question then becomes: how can these be supported when repeated interactions are not guaranteed? Requires some formation of a partial tendency to not defect in order to support broader social cooperation -- i.e., some initial impetus that can then be built upon
\end{itemize}


\section{Stinchcombe: ``Reason and Rationality''}

\begin{itemize}
    \item Specifies an ideal long-run relationship between Weber's substantive and formal rationality:
    \begin{itemize}
        \item Formal rationality: ``standardized methods of calculation on which routines can be based'' -- strict and explicit institutions
        \item Substantive rationality: Beyond formal methods and instead engaging with substance -- how strong is your CV, is monetary gain a sufficient incentive, etc.
    \end{itemize}
    \item ``What we want of our institutions of reason is a proper balance between efficient formal approximations that can have a reliable social effect, and substantive good sense to know their limits and to improve them.''
\end{itemize}


\section{Miller: ``Managerial Dilemmas: Political Leadership in Hierarchies''}

\begin{itemize}
    \item Two distinct literatures on organizational control:
    \begin{enumerate}
        \item Control by means of designing aligned incentive systems and sanctions such that self-interested / utility-maximizing actors behave in some desired way
        \item Control by means of behavioral influence -- inspiration, social encouragement, generally means that are not tied to some form of explicit payoff or sanction
    \end{enumerate}
    \item General claim: this first approach is self-destructive, and is demonstrated by social choice theory to be logically impossible. The most effective approach to control is subsequently personalistic leadership (tie to Weber)
    \item ``Once the Leviathan is created to correct the deficiencies of collective action, Leviathan itself is found to generate inefficiencies. The very mechanism that makes Leviathan desirable -- the authority to change members' incentives -- is found to be logically insufficient to guarantee efficiency. The irony, then, is that the managers of a hierarchy will find it desirable to look for modes of inducing voluntary cooperation within the Leviathan.'
\end{itemize}


\section{Hardin: The Social Evolution of Cooperation}

\begin{itemize}
    \item Coordination as a foundational principle for cooperation -- we must be able to coordinate on some cooperative process, if we are to achieve (and sustain) that cooperation
    \item Gunman theory of law: ``if we are to be made to obey law by threat of force, then the state will be unable to mount adequate mechanisms of enforcement.''
    \item Perverse logic of state construction: if we construct the state to enforce collective action solutions, then we must overcome the collective action problem of constructing the state itself
    \item Common law as a reflection of path-dependency
    \item Coordination creates a convention, and a convention promotes further coordination -- circular pathway
\end{itemize}


\section{North: ``Institutions and Their Consequences for Economic Performance''}

\begin{itemize}
    \item Largely consistent with North 1990a -- institutions reduce transaction costs
    \item ``Institutions consist of informal constraints and formal rules and of their enforcement characteristics.''
    \item Effectiveness of institutions is largely reflected through how much they reduce transaction costs
    \item Cites North 1981 -- two reasons for institutions producing lesser property rights. First, may pander to certain powerful groups who may otherwise support an alternative ruler/regime. Second, may seek to gain tax revenue (which could be decreased through stronger property rights)
    \item North 1990b: ``a basic function of institutions is to provide stability and continuity by dampening the effects of relative price changes. It is institutional stability that makes possible complex exchange across space and time.''
\end{itemize}


\section{Levi: ``A Logic of Institutional Change''}




\end{document}
